\documentclass[aspectratio=169, 10pt]{beamer}
\usepackage[sfdefault]{roboto}
\usepackage{roboto-mono}
\usepackage[utf8]{inputenc}
\usepackage[T1]{fontenc}
\usepackage[portuguese]{babel}
\usepackage{booktabs}
\usepackage{adjustbox}
\usepackage{color}
\usepackage{amsmath}
\usepackage{amssymb}
\usepackage{minted}
\usepackage{tikz}
\usetikzlibrary{shapes.geometric, arrows.meta, positioning, calc}
\tikzstyle{startstop} = [rectangle, rounded corners, minimum width=3cm, minimum height=1cm, text centered, draw=black, fill=red!20]
\tikzstyle{process} = [rectangle, minimum width=3cm, minimum height=1cm, text centered, draw=black, fill=orange!20]
\tikzstyle{decision} = [diamond, minimum width=3cm, minimum height=1cm, text centered, draw=black, fill=green!20]
\tikzstyle{arrow} = [thick, ->, >=Stealth]

% Theme configuration
\usetheme{Madrid}
\usecolortheme{whale}

\setbeamertemplate{itemize item}{\small\color{brightblue}$\blacktriangleright$}
\setbeamertemplate{itemize subitem}{\tiny\color{mainblue}$\bullet$}

\definecolor{mainblue}{HTML}{2a3990}
\definecolor{brightblue}{HTML}{3c78d8}
\definecolor{lightblue}{HTML}{ebf0f9}
\definecolor{darkgreen}{HTML}{1e7e34}
\definecolor{lightgreen}{HTML}{d4edda}
\definecolor{darkred}{HTML}{bd2130}
\definecolor{lightred}{HTML}{f8d7da}
\definecolor{darkorange}{HTML}{d97706}
\definecolor{lightorange}{HTML}{fef3c7}

\setbeamercolor{frametitle}{fg=mainblue, bg=white}
\setbeamerfont{frametitle}{series=\bfseries, size=\Large}

\setbeamertemplate{frametitle}{
    \begin{beamercolorbox}[wd=\paperwidth, leftskip=1.2em, center, rightskip=1.2cm, ht=2.0em, dp=0.8em]{frametitle}
        \centering\usebeamerfont{frametitle}\insertframetitle
    \end{beamercolorbox}
}

\setbeamersize{text margin left=2em, text margin right=2em}
\setlength{\parindent}{0pt}

\beamertemplatenavigationsymbolsempty
\setbeamertemplate{navigation symbols}{}
\setbeamertemplate{footline}{
  \hfill%
  \usebeamerfont{page number in head/foot}%
  \usebeamercolor[fg]{page number in head/foot}%
  \small\insertframenumber%
  \hspace*{2em}\vspace*{1.2em}%
}

\title[ARC-AGI: CEGIS e Anti-Overfitting]{S\'{\i}ntese de Programas com CEGIS e Restri\c{c}\~{o}es Anti-Overfitting no ARC-AGI}
\author{Fabio Luiz da Costa Cieslak \and Izadora Candotti de Oliveira \and Matheus Machado}
\institute{Universidade Federal do Rio Grande do Sul (UFRGS)}
\date{2026}

\begin{document}

% -------------------------------------------------------------
% SLIDE 1: Capa
% -------------------------------------------------------------
\begin{frame}[plain]
  \centering
  \color{mainblue}
  
  {\huge \bfseries S\'{\i}ntese de Programas com CEGIS e Restri\c{c}\~{o}es Anti-Overfitting no ARC-AGI}\\[0.5em]

  \smallskip
  \noindent\rule{\textwidth}{1.5pt}
  \medskip

  \textbf{Fabio Luiz da Costa Cieslakk, Izadora Candotti de Oliveira e Matheus Machado}\\[0.3em]
  {\small Universidade Federal do Rio Grande do Sul}\\[1cm]
\end{frame}

% -------------------------------------------------------------
% SLIDE 2: Contexto: ARC-AGI e Síntese de Programas
% -------------------------------------------------------------
\begin{frame}[fragile]{ARC-AGI}
    \vspace{1.0em}
    
    \begin{center}
        \Large Desafio de \textbf{racioc\'{\i}nio indutivo} baseado em transformações de grids coloridos a partir de demonstra\c{c}\~{o}es visuais
    \end{center}

    
\end{frame}

\begin{frame}[fragile]{S\'{\i}ntese de Programas}
    \vspace{1.0em}
    \begin{itemize}\setlength{\itemsep}{1.1em}
        \item S\'{\i}ntese de fun\c{c}\~{a}o em Python que busca resolver a tarefa de transformação:
            \begin{center}
                \begin{minted}[fontsize=\small, frame=none]{python}
def transform(grid: list[list[int]]) -> list[list[int]]:
                \end{minted}
            \end{center}
        \item Criado a partir da grade de \textbf{treinamento}
        \item Avalia\c{c}\~{a}o por acerto bin\'{a}rio na grade de teste \textbf{oculta}
    \end{itemize}
\end{frame}
% -------------------------------------------------------------
% SLIDE 3: Abordagens Avaliadas: Baseline vs. CEGIS
% -------------------------------------------------------------
\begin{frame}{Abordagens Avaliadas}
    \vspace{1.0em}
    \textbf{Baseline}:\vspace{0.5em}
    \begin{itemize}\setlength{\itemsep}{1.1em}
        \item Gera\c{c}\~{a}o est\'{a}tica em \textbf{tentativa \'{u}nica (1-Shot)}, sem mecanismos de reparo
    \end{itemize}
    \vspace{1.0em}
    \textbf{CEGIS}:\vspace{0.5em}
    \begin{itemize}\setlength{\itemsep}{1.1em}
        \item Refinamento \textbf{iterativo guiado por contraexemplos (CEGIS)} em at\'{e} 5 rodadas
        \item Retorno do \textbf{primeiro contraexemplo} \`{a} LLM para corre\c{c}\~{a}o de c\'{o}digo
    \end{itemize}
\end{frame}

% -------------------------------------------------------------
% SLIDE 4: Estratégia Proposta (AntiCheat) e Hipótese
% -------------------------------------------------------------
\begin{frame}{Estrat\'{e}gia Proposta (AntiCheat)}
    \vspace{0.6em}
    \begin{itemize}\setlength{\itemsep}{0.9em}
        \item Proibi\c{c}\~{a}o de uso de \textbf{atalhos} ou \textbf{coordenadas fixas} nas grades de treino
        \item Obrigatoriedade de enunciar o \textbf{invariante abstrato} antes de emitir o c\'{o}digo
    \end{itemize}
    \vspace{0.8em}
    \begin{center}
        {\color{mainblue}\rule{0.8\linewidth}{0.5pt}}\\[0.6em]
        \textbf{Hip\'{o}tese Central}
    \end{center}
    \begin{quote}
        \small\centering\itshape
        ``Restri\c{c}\~{o}es sem\^{a}nticas anti-hardcoding e reflex\~{a}o sobre invariantes melhoram a acur\'{a}cia no teste oculto do ARC-AGI''
    \end{quote}
\end{frame}

% -------------------------------------------------------------
% SLIDE 5: Escala Experimental e Protocolo Pareado
% -------------------------------------------------------------
\begin{frame}{Escala Experimental}
    \vspace{0.8em}
    \begin{itemize}\setlength{\itemsep}{0.9em}
        \item Amostra de \textbf{100 tarefas} avaliadas em \textbf{8 modelos} (\textbf{800 testes pareados})
        \item Cobertura em tr\^{e}s faixas de capacidade de LLMs:
            \begin{itemize}\setlength{\itemsep}{0.3em}
                \item \textbf{Fronteira} com Grok 4.6, DeepSeek V4 Flash e GPT-5.6 Luna
                \item \textbf{Intermedi\'{a}rios} com Claude Sonnet 5 e fam\'{\i}lia Gemini Flash
                \item \textbf{Compactos} com Leanstral 1.5.1 e GPT-5.4 Nano
            \end{itemize}
    \end{itemize}
\end{frame}

% -------------------------------------------------------------
% SLIDE 7: Eficácia do CEGIS sobre o Baseline (TABELA SOZINHA)
% -------------------------------------------------------------
\begin{frame}{Efic\'{a}cia do CEGIS sobre o Baseline (1-Shot)}
    \centering
    \vspace{0.5em}
    \small
    \begin{tabular}{lccc}
        \toprule
        \textbf{Modelo} & \textbf{Baseline (1-Shot)} & \textbf{CEGIS (Iterativo)} & \textbf{Ganho Absoluto ($\Delta$)} \\
        \midrule
        Grok 4.6 & 88\% & 92\% & +4\% \\
        DeepSeek V4 Flash 0731 & 74\% & 86\% & +12\% \\
        GPT-5.6 Luna & 58\% & 77\% & +19\% \\
        Claude Sonnet 5 & 32\% & 62\% & +30\% \\
        Gemini 3.5 Flash Lite & 21\% & 43\% & +22\% \\
        Gemini 3.1 Flash Lite & 12\% & 31\% & +19\% \\
        Leanstral 1.5.1 & 3\% & 7\% & +4\% \\
        GPT-5.4 Nano & 1\% & 5\% & +4\% \\
        \midrule
        \textbf{Total Agregado (800 testes)} & \textbf{36,1\%} & \textbf{50,4\%} & \textbf{+14,3\%} \\
        \bottomrule
    \end{tabular}
\end{frame}

% -------------------------------------------------------------
% SLIDE 8: Resultados Globais: CEGIS vs. AntiCheat (TABELA SOZINHA)
% -------------------------------------------------------------
\begin{frame}{Resultados Globais: CEGIS vs. AntiCheat}
    \centering
    \vspace{0.5em}
    \footnotesize
    \begin{tabular}{lccccc}
        \toprule
        \textbf{Modelo} & \textbf{CEGIS (\%)} & \textbf{AntiCheat (\%)} & \textbf{Vit\'{o}rias ($b / c$)} & \textbf{$p$-valor bruto} & \textbf{$p$ ajustado (Bonf.)} \\
        \midrule
        Grok 4.6 & 92,0\% & 92,0\% & 0 / 0 & 1,0000 & 1,0000 \\
        DeepSeek V4 Flash 0731 & 86,0\% & 86,0\% & 0 / 0 & 1,0000 & 1,0000 \\
        GPT-5.6 Luna & 77,0\% & 79,0\% & 2 / 4 & 0,6875 & 1,0000 \\
        Claude Sonnet 5 & 62,0\% & 60,0\% & 5 / 3 & 0,7266 & 1,0000 \\
        Gemini 3.5 Flash Lite & 43,0\% & 45,0\% & 6 / 8 & 0,7905 & 1,0000 \\
        Gemini 3.1 Flash Lite & 31,0\% & 24,0\% & 7 / 0 & 0,0156 & 0,1250 \\
        Leanstral 1.5.1 & 7,0\% & 9,0\% & 2 / 4 & 0,6875 & 1,0000 \\
        GPT-5.4 Nano & 5,0\% & 3,0\% & 3 / 1 & 0,6250 & 1,0000 \\
        \midrule
        \textbf{Total Agregado} & \textbf{50,4\%} & \textbf{49,8\%} & \textbf{25 / 20} & \textbf{0,5515} & --- \\
        \bottomrule
    \end{tabular}
\end{frame}

% -------------------------------------------------------------
% SLIDE 6: Distribuição de Resolução das Tarefas (IMAGEM SOZINHA)
% -------------------------------------------------------------
\begin{frame}{Distribui\c{c}\~{a}o de Resolu\c{c}\~{a}o das Tarefas}
    \centering
    \includegraphics[width=\linewidth, height=0.79\textheight, keepaspectratio]{figures/plot_task_solve_distribution.png}
\end{frame}

% -------------------------------------------------------------
% SLIDE 9: Matriz de Discrepância das Estratégias (IMAGEM SOZINHA)
% -------------------------------------------------------------
\begin{frame}{Matriz de Discrep\^{a}ncia das Estrat\'{e}gias}
    \centering
    \includegraphics[width=\linewidth, height=0.79\textheight, keepaspectratio]{figures/plot_strategy_discrepancy.png}
\end{frame}

% -------------------------------------------------------------
% SLIDE 10: Testes de Significância Estatística (TABELA SOZINHA)
% -------------------------------------------------------------
\begin{frame}{Testes de Signific\^{a}ncia Estat\'{\i}stica}
    \centering
    \vspace{0.8em}
    \resizebox{0.95\linewidth}{!}{%
    \begin{tabular}{lccc}
        \toprule
        \textbf{Teste Estat\'{\i}stico} & \textbf{Valor / Intervalo} & \textbf{Conclus\~{a}o} \\
        \midrule
        Teste Binomial Exato & $p = 0{,}5515$ & N\~{a}o significante \\
        Teste de Permuta\c{c}\~{a}o & $p = 0{,}5501$ & N\~{a}o significante \\
        Intervalo de Confian\c{c}a Bootstrap & $[-2{,}25\%, +1{,}00\%]$ & Cont\'{e}m o zero \\
        Corre\c{c}\~{a}o de Bonferroni ao Gemini 3.1 & $p_{\text{adj}} = 0{,}1250$ & N\~{a}o significante \\
        \bottomrule
    \end{tabular}%
    }
\end{frame}

% -------------------------------------------------------------
% SLIDE 12: Diagnóstico Qualitativo: Viés e Gargalo Real
% -------------------------------------------------------------
\begin{frame}{Diagn\'{o}stico Qualitativo: Vi\'{e}s e Gargalo Real}
    \vspace{1.0em}
    \begin{itemize}\setlength{\itemsep}{1.1em}
        \item C\'{o}digos gerados sob CEGIS e AntiCheat s\~{a}o \textbf{estruturalmente quase id\^{e}nticos}
        \item Mais de \textbf{$90\%$ das falhas} decorrem da incapacidade de resolver os treinos (\textbf{teto de representa\c{c}\~{a}o})
    \end{itemize}
\end{frame}

% -------------------------------------------------------------
% SLIDE 13: Métricas Psicométricas de Guttman (TABELA SOZINHA)
% -------------------------------------------------------------
\begin{frame}{Hierarquia de Tarefas}
    \vspace{1.0em}
    \begin{itemize}\setlength{\itemsep}{1.1em}
        \item Tarefas seguem um modelo quase que totalmente \textbf{hierárquico}
        \item Grok 4.6 (1º) acertou \textbf{100\%} das tarefas que o DeepSeek V4 (2º) acertou
        \item DeepSeek V4 (2º) acertou \textbf{97,4\%} das tarefas que o GPT-5.6 Luna (3º) acertou
    \end{itemize}
\end{frame}

% -------------------------------------------------------------
% SLIDE 16: Dinâmica de Recuperação por Iteração (IMAGEM SOZINHA)
% -------------------------------------------------------------
\begin{frame}{Din\^{a}mica de Recupera\c{c}\~{a}o por Itera\c{c}\~{a}o}
    \centering
    \includegraphics[width=\linewidth, height=0.79\textheight, keepaspectratio]{figures/plot_iteration_diminishing_returns.png}
\end{frame}

% -------------------------------------------------------------
% SLIDE 17: Dinâmica de Iterações do CEGIS (TABELA SOZINHA)
% -------------------------------------------------------------
\begin{frame}{Din\^{a}mica de Itera\c{c}\~{o}es do CEGIS}
    \centering
    \vspace{0.5em}
    \small
    \begin{tabular}{ccccc}
        \toprule
        \textbf{Itera\c{c}\~{a}o} & \textbf{Tarefas Resolvidas} & \textbf{\% do Total Resolvido} & \textbf{Tarefas Recuperadas} & \textbf{\% Acumulado} \\
        \midrule
        1 & 289 & 71,7\% & 0 & 0,0\% \\
        2 & 69 & 17,1\% & 69 & 60,5\% \\
        3 & 29 & 7,2\% & 29 & \textbf{86,0\%} \\
        4 & 9 & 2,2\% & 9 & 93,9\% \\
        5 & 7 & 1,7\% & 7 & 100,0\% \\
        \bottomrule
    \end{tabular}
\end{frame}
% -------------------------------------------------------------
% SLIDE 19: Conclusões
% -------------------------------------------------------------
\begin{frame}{Conclus\~{o}es}
    \vspace{1.0em}
    \begin{itemize}\setlength{\itemsep}{1.1em}
        \item Refinamento por \textbf{CEGIS supera o Baseline em $+14{,}3\%$} (114 tarefas recuperadas)
        \item Restri\c{c}\~{o}es \textbf{AntiCheat s\~{a}o estatisticamente equivalentes} ao CEGIS padr\~{a}o ($p = 0{,}5515$)
        \item Principal gargalo reside no \textbf{teto de representa\c{c}\~{a}o ($>90\%$)}, e n\~{a}o em trapa\c{c}as ($0{,}29\%$)
    \end{itemize}
\end{frame}

% -------------------------------------------------------------
% SLIDE 20: Encerramento
% -------------------------------------------------------------
\begin{frame}{Encerramento}
    \centering
    \vspace{2em}
    {\Huge\color{mainblue}\textbf{Obrigado!}}\\[1em]
    {\large\color{gray}Perguntas e Discuss\~{a}o}\\[2em]
    \normalsize
    Grupo C -- UFRGS\\
    F\'{a}bio Cieslak $\bullet$ Izadora Candotti $\bullet$ Matheus Machado
\end{frame}

\end{document}
